% ===================================================================
% Paper 4, v5.5 (2026-07)
% Dark energy as the instanton residual of the condensate's angular
% mode. Solve-oriented A+B construction: scheme-separated floor +
% surviving field-dependent residual; dimensional-transmutation scale;
% thawing PNGB quintessence; unification with the galactic
% acceleration scale a0 ~ Lambda_DE^2 / M_Pl (MOND-Lambda coincidence).
% v5.3: central computation re-scoped (fermion-kernel index excluded).
% v5.5: single-condensate decay-constant fix; third chi_top candidate
% (eta-Hopf, Paper 5 S2-3) registered; version labels synced.
% v5.6: (i) mediator-gap section added (parity theorem + two computed
% negative results close the loop-level routes; topological funnel);
% (ii) smooth-field Hopf density proven identically zero (pullback
% theorem, machine-sealed); eta-Hopf candidate re-scoped to the
% MONOPOLE sector of the orientation layer; (iii) monopole viability
% gate PASSED: measured Skyrme band brackets the required loop action
% S ~ 276; S2-2 co-resolved (Coulomb phase); (iv) f/M_Pl thawing link
% recorded; (v) erratum: b0 g^2 ~ 0.29 (not ~1) for S = 276.
% v5.7: MAGI round 1 (ChatGPT) corrections: (i) 3+1D vs 2+1D
% compact-U(1) physics fixed -- Coulomb phase has an EXACTLY massless
% photon (Guth); "Polyakov-class mass" phrasing removed; correction
% propagates to Paper 5 gate S2-2 (restated as a phase question,
% answered on the Coulomb side); (ii) gate claim scoped: monopole
% DILUTENESS scale passes; chi_top itself awaits linked/twisted-loop
% computation (C2b).
% v5.8: abstract trimmed to arXiv limit (1888 chars, was 3669); the
% displaced technical statements already live in Secs. 2, 7, 8.
% Review process: per-version review = Claude + ChatGPT; fresh-session
% re-review before submission.
% v6.0: MAJOR REVISION following the C2a measurement.
% (1) C2a executed: first chirally complete kappa_U measurement
% (pointwise A0=0 to machine precision, five diagram classes, two
% evaluators, two integrators): kappa_U = -0.00149(1) -- 97x below the
% required 0.144 and of opposite sign. S_mono = -0.95, outside
% [140,550]: THE TRANSMUTATION/MONOPOLE CHAIN TERMINATES by its
% pre-registered criterion. v5.6-5.8 viability-gate claims (beta band,
% Coulomb co-resolution) RETRACTED -- they rested on a fit-leakage
% artifact in the July-8 kappa numbers.
% (2) Ontology repositioned on the author's ground: internal-dimension
% energy reservoir ('the sea'), lattice vacuum as its surface.
% (3) New candidate answer for CC smallness: Volovik equilibrium
% argument (self-sustained condensate => zero vacuum pressure; only
% deviations gravitate), with its known weakness stated; NEW CENTRAL
% COMPUTATION registered: audit of the induced Einstein equations'
% source term in the companion gravity paper's machinery.
% (4) Outlook (picture, not derivation): expansion as lattice
% accretion; constant rho_DE as bookkeeping identity; falsifiable
% edge at bound systems.
% (5) Thawing phenomenology downgraded to conditional (potential
% source unidentified; RtildeR route sole registered backup).
% v6.3: cross-review sweep (ChatGPT, three-paper round): DESI/thawing
% passage made strictly conditional and equilibrium-vacuum's
% no-commitment status stated; Summary bullet 1 rewritten (candidate
% carrier, amplitude unexplained, active explanation = equilibrium
% subtraction); galactic-unification bullet retired as historical
% motivation with provenance; Acknowledgments C2a-C2d status
% corrected to executed.
% v6.2: sec 8.4 updated -- omega existence+dynamization gates passed
% (companion note); stabilization consequence now has an in-framework
% candidate; photon consequence unresolved; R_star gate open.
% v6.1: narrative-reset sweep (MAGI round, ChatGPT): all residual
% v5.x claims removed -- 'scale is generated' -> generic possibility
% with explicit no-computed-source status; 'dark-energy axion' ->
% 'candidate'; 'smallness follows from transmutation' -> ledger
% rewrite (form+smoothness survive; kind conditional; amplitude
% unexplained); 'mechanism predicts thawing' -> conditional PNGB;
% intro thesis rewritten as falsification-and-repositioning; status
% box added at eq:transmute. Paper's stated identity: no longer a
% derived-mechanism paper.
% (6) Author-comprehension (teach-back) gate instituted: all theory
% content in this version restated and approved by the author.
% Companion to: Paper 1 (SPARC / emergent a0, v12), Paper 2 (emergent
% GR, v2.14), Paper 3 (vector channel, v3.10), Paper 5 (working note,
% v0.19).
% ===================================================================
\documentclass[11pt]{article}
\usepackage{amsmath,amssymb,graphicx,booktabs,geometry,xcolor,slashed}
\geometry{margin=1in}
\newcommand{\beq}{\begin{equation}}
\newcommand{\eeq}{\end{equation}}
\newcommand{\Mpl}{M_{\mathrm{Pl}}}
\newcommand{\rhoL}{\rho_\Lambda}
\newcommand{\mchi}{m_\chi}
\newcommand{\LDE}{\Lambda_{\mathrm{DE}}}

\title{Dark Energy as the Instanton Residual of a Lattice Fermion
Condensate:\\
Pseudo-Goldstone Quintessence and the Origin of the
MOND--$\Lambda$ Coincidence}
\author{Z.~Cheng}
\date{July 2026 --- v6.3 (draft)}

\begin{document}
\maketitle

\begin{abstract}
We reposition the dark-energy sector of the lattice fermion condensate
framework on measured ground.  The ontology: the condensate's
internal-dimension sector is an energy reservoir---a deep sea---of
which the four-dimensional lattice vacuum is the surface; parity
protects the smoothness of the surface's angular mode, and only
structure (waves, knots, defects) should carry weight.
The question of why the observed dark energy sits $122$ orders below
the Planck scale is then addressed at three levels.
(i) A candidate answer of Volovik type: a self-sustained condensate in
equilibrium has identically vanishing vacuum pressure, so the
equilibrium sea does not gravitate; whether the framework's
\emph{induced} gravity realizes this---whether the emergent Einstein
equations source only deviations from equilibrium---is a well-posed
audit in the companion gravity paper's machinery, registered here as
the new central computation.
(ii) An executed and terminated mechanism: the dimensional-transmutation
route ($\LDE^4 \sim \Lambda^4 e^{-S}$ via dilute monopoles of the
orientation layer) required a lattice Skyrme coefficient
$\kappa_U \simeq 0.14$; the first chirally complete measurement---five
diagram classes, pointwise $A_0 \equiv 0$ to machine precision, two
independent evaluators and integrators---gives
$\kappa_U = -0.00149 \pm 0.00001$: ninety-seven times short and of
opposite sign.  $S_{\rm mono} = -0.95 \notin [140,550]$, and the chain
terminates by its pre-registered criterion; the viability-gate claims
of earlier versions, which rested on a fit-leakage artifact, are
retracted.
(iii) An outlook, marked as picture: expansion as lattice
accretion, with constant $\rho_{\rm DE}$ as a bookkeeping identity
and a falsifiable edge at bound systems.
What survives measurement: the parity smoothness theorem, the
smooth-field Hopf-density theorem, the honest cosmological-constant
accounting, and the unexplained scale coincidence
$a_0 \sim \LDE^2/\Mpl$.
\end{abstract}

% ===================================================================
\section{Introduction}
\label{sec:intro}
% ===================================================================

The lattice fermion condensate programme treats spacetime, gravity,
and dark matter as collective phenomena of a Planck-scale fermion
substrate.  The gravitational sector is developed in
\cite{Cheng:2026grav}, the galactic sector (a single emergent acceleration scale $a_0$,
established against SPARC rotation curves after the decisive exclusion
of the earlier fixed-range Yukawa kernel) in
\cite{Cheng:2025sparc}, and the vector channel that controls the
sign of the induced Newton constant in \cite{Cheng:2026vector}.
Dark energy was the subject of the earliest version of the vector
paper, under the title \emph{Dark Energy from a Pseudo-Hermitian
Vacuum}; that mechanism has since been retired
\cite{Cheng:2026vector}.

This paper develops the dark-energy sector of that programme.  Its
thesis has been reshaped by measurement, and this version states the
reshaped form: dark energy is carried by the condensate's angular
mode---a pseudo-Goldstone whose \emph{smoothness} is protected by
parity (a theorem that survives everything below)---while the
\emph{amplitude} $\LDE^4$ is at present unexplained: the concrete
mechanism this programme constructed for it (dimensional transmutation
through the orientation layer's monopole gas) was pre-registered,
executed, and terminated by its own criterion
(Section~\ref{sec:monopole}).  The paper is accordingly no longer a
derived-mechanism paper; it is a falsification-and-repositioning
paper, whose open front is the equilibrium-vacuum audit of
Section~\ref{sec:volovik}.  We are
equally explicit about the one piece the construction does \emph{not}
supply---why the cancelled vacuum floor leaves a near-zero residual,
the standard cosmological-constant problem shared with every
induced-gravity programme---and we adopt the same posture as the
companion gravitational paper \cite{Cheng:2026grav}, recording open
problems sharply rather than concealing them.  What is new here is that
the residual cosmological constant, once separated as a counterterm
convention, no longer gates the \emph{dynamical} dark energy: the
field-dependent instanton term survives the counterterm and is the
physical object we compute.

Throughout, $\rhoL \equiv \rho_{\mathrm{DE}}$ denotes the observed
dark-energy density, $\rhoL \approx (2.3\times10^{-3}\,
\mathrm{eV})^4 \approx 2.8\times10^{-11}\,\mathrm{eV}^4$, and
$\Mpl = 2.4\times10^{27}\,\mathrm{eV}$ is the reduced Planck mass.

% ===================================================================
\section{The dark-energy axion candidate}
\label{sec:closed}
% ===================================================================

\subsection{A retired alternative: the imaginary vacuum}
\label{sec:imaginary}

The complex condensate is $\Phi = \rho\,e^{i\theta}$, and the
substrate fermion couples to it through a \emph{chiral} mass
$M = g(\sigma + i\gamma_5\pi) = g\rho\,e^{i\gamma_5\theta}$, in which
the pseudoscalar component necessarily carries $\gamma_5$.  The
Dirac determinant of such a mass depends only on the modulus,
\beq
\det\bigl(i\slashed{p}_E + g\rho\,e^{i\gamma_5\theta}\bigr)
= \bigl(p_E^2 + g^2\rho^2\bigr)^2 ,
\eeq
independent of $\theta$, because the chiral rotation assigns
opposite phases to the left- and right-handed components, which
cancel pairwise in the closed loop.  The one-loop angular direction
is therefore flat, and every value of $\theta$---including
$\theta = \pi/2$, the Hermitian pseudoscalar configuration
$g\rho\,i\gamma_5$---is Hermitian.

The earlier dark-energy mechanism rested on a purely imaginary,
non-Hermitian mass $i g\rho$ (``$\theta = \pi/2$''), held metastable
by a $\gamma_5$-pseudo-Hermitian, Krein-space construction, with the
cosmological constant identified as the energy difference $\Delta V$
between this configuration and the real vacuum.  Two independent
facts retire it \cite{Cheng:2026grav,Cheng:2026vector}.  The
configuration $ig\rho$ is not a point of the chiral condensate
manifold at all, so the ``$\theta = \pi/2$ maximum'' of the older
$\cos2\theta$ landscape was an artifact of a non-chiral
parametrization of the determinant.  And any non-Hermitian
construction fails on its own terms: the dispersion relation at such
a configuration has imaginary energies for $|p| < m_I$, so the
relevant $\mathcal{PT}$ symmetry is broken in the infrared and the
reality theorems do not apply.  The intuition that ``energy stored
on the imaginary axis drives expansion'' does not survive contact
with the corrected vacuum structure.

\subsection{The angular mode as dark-energy axion candidate}

\textit{Structural requirement (added in v5.1).}  The
identification below requires the microscopic four-fermion
structure to preserve the classical $U(1)_A$, so that the angular
mode's potential arises \emph{only} from the topological sector;
tree-level $U(1)_A$ breaking would lift the mode to the cutoff
scale.  In the multi-flavor extension of the programme this
selects the classically $U(2)\times U(2)$-symmetric interaction
(machine-verified to preserve the vacuum-stability protections;
see the internal-dimension working note), whose vacuum manifold
$U(2) \cong (S^1 \times S^3)/\mathbb{Z}_2$ carries this paper's
angular mode on the $S^1$ factor and topological matter on the
$S^3$ factor.
The completed re-audit adds three quantitative statements.
First, the classical $U(1)_A$ invariance of the
$U(2)\times U(2)$ interaction is \emph{exact} (machine-verified
under continuous axial rotations), so the tree-level flatness of
the angular direction is a symmetry theorem rather than an
approximation.
Second, and decisively for methodology: in the Wilson
discretization used for the programme's lattice campaigns, the
mechanism of this paper does not merely suffer contamination---it
is \emph{absent}: at the condensate operating point the
pseudoscalar channels are bubble-supercritical
($G_{\rm op} = 2.97$ against a pseudoscalar bubble criticality of
$1.51$ in the same calibrated units), so the would-be Goldstone
directions are reorganized at the cutoff scale (Aoki-type).
A Ginsparg--Wilson (overlap) discretization is therefore a
\emph{requirement} of the dark-energy sector, not an improvement:
it preserves the exact tree-level flatness while supplying the
anomaly through the fermion measure, and its index theorem gives
the first well-posed lattice definition of the topological charge
whose susceptibility is this paper's central open computation.
Third, minor bookkeeping: the flavor trace doubles the angular
mode's kinetic normalization, $f \to \sqrt{2} f$, slightly
easing the super-Planckian decay-constant requirement, and the
$U(2)$ model adds $\pi_4(SU(2)) = \mathbb{Z}_2$ to the candidate
topological structures relevant to that computation.
\textit{Re-scoping of the central computation (v5.3).}
A machine test closes one candidate and re-scopes the search: with
trivial links, the overlap topological charge
$Q = -\tfrac12\,{\rm Tr}\,{\rm sign}(H_W)$ vanishes identically
for every chiral-field background tested (uniform $U(1)_A$
windings and localized lumps), confirming the analytic statement
that a pure fermion substrate carries no $U(1)_A$ anomaly: the
angular mode is an \emph{exact} Goldstone at this level, and the
potential this paper requires cannot arise from the fermion-kernel
index.  A $F\tilde F$-like source is needed.  Within the
programme three candidates remain, all uncomputed: a composite
gauge sector (disfavored at natural couplings by the
vector-channel pole analysis of the companion paper); the
mixed $U(1)_A$--gravitational anomaly of the \emph{emergent
geometry itself}, $\partial_\mu j_A^\mu \propto R\tilde R$, under
which $\Lambda_{\rm DE}$ would be set by the topological
susceptibility of emergent-geometry fluctuations---a speculative
but structurally natural route that would make the dark
energy--gravity connection intrinsic rather than coincidental;
and \emph{(added in v5.5)} the Hopf density $f \wedge f$ of the
composite Berry connection carried by the $\mathbb{CP}^1$
orientation quotient of the $U(2)$ chiral field (the $S^2$ layer
of the internal-dimension working note, gate S2-3): whether the
fermion determinant couples the angular mode to this density is a
well-posed derivative-expansion question, and a positive answer
would make Hopfion topological fluctuations of the orientation
sector the source, at no cost in new degrees of freedom.
The central open computation of this paper is accordingly
re-scoped from ``identify the topological sector'' to ``compute
the gravitational topological susceptibility of the emergent
geometry, exhibit a composite gauge structure, or establish the
$\eta$--Hopf coupling of the orientation sector.''
\textit{Second re-scoping (v5.6).}
The third candidate is now sharpened by a theorem.
The composite Berry curvature of the $\mathbb{CP}^1$ orientation
layer is the pullback of the Fubini--Study two-form,
$f = \phi^*\omega$; since $\omega\wedge\omega = 0$ on a
two-manifold and pullback commutes with the wedge product,
$f \wedge f \equiv 0$ \emph{pointwise for every smooth
configuration} (machine-verified:
$|\varepsilon^{\mu\nu\rho\sigma}f_{\mu\nu}f_{\rho\sigma}|/|f|^2
< 5\times10^{-16}$ over $200$ random smooth fields).
The Hopf-density candidate as originally stated is therefore empty
on smooth fields, and the susceptibility can arise only from the
\emph{singular} sector: the $\pi_2(S^2) = \mathbb{Z}$ monopole
defects whose worldline loops carry non-exact flux---precisely the
objects that gate S2-2 of the internal-dimension working note had
flagged as the photon candidate's principal assassin.
Sections~\ref{sec:mediator} and~\ref{sec:monopole} develop the
consequences: the composite gauge route remains disfavored at
natural couplings \cite{Cheng:2026vector}, the gravitational
$R\tilde R$ route remains open as a backup, and the monopole sector
becomes the leading candidate, with a passed viability gate.

\label{sec:axion}

The phase of the condensate is fixed not at one loop (which is flat in
the corrected chiral parametrization) but by the subleading
chiral-symmetry-breaking effects---the axial anomaly and the discrete
lattice structure---which generate a shallow periodic potential of
instanton origin,
\beq
V_\theta(\theta) = K\,[\,1 - \cos(M\theta)\,],
\qquad K \sim \LDE^4 ,
\label{eq:Vtheta}
\eeq
with minimum at the $CP$-conserving point $\theta = 0$.  We identify
this angular pseudo-Goldstone as the dark-energy field.  Two features
make the identification natural and were not available to the earlier,
imaginary-vacuum reading.

\emph{It is automatically smooth.}  The angular mode is a pseudoscalar;
parity forbids a non-derivative coupling to the scalar density of
unpolarized matter, $\phi\,\bar\psi\psi$, so the field does not cluster
and mediates no fifth force---exactly the inertness dark energy
requires, obtained here for free rather than imposed.  (The same
parity argument places the matter-coupled, fifth-force-bearing degree
of freedom in the \emph{scalar} channel, which is where the companion
galactic paper locates it; see Section~\ref{sec:mond}.)

\emph{Status (v6.0): no computed sector currently supplies the
scale.}  The height $K\sim\LDE^4$ would be a topological
susceptibility, and a susceptibility \emph{could in principle} lie
exponentially below the cutoff by dimensional transmutation,
\beq
\LDE^4 \;\sim\; \Lambda^4\, e^{-S}, \qquad
S \;=\; \frac{8\pi^2}{b_0\, g^2(\Lambda)} ,
\label{eq:transmute}
\eeq
so that $\rhoL/\Mpl^4 \sim 10^{-120}$ follows from
$b_0\,g^2(\Lambda)\simeq 0.29$ (i.e.\ $S \simeq 276$)---an
$\mathcal{O}(1)$ coupling, e.g.\ $g \simeq 0.5$ with
$b_0 \simeq 1.1$---rather than a $122$-digit cancellation.
(Erratum for v5.5, which stated $b_0 g^2 \simeq 1$; that value
would give $e^{-S} \sim 10^{-34}$, not $10^{-120}$.)
This is a \emph{generic possibility, not a result of the framework}:
the one concrete realization this programme constructed---the
monopole gas of the orientation layer, with
$1/g'^2 = 32N\kappa_U$---fails the pre-registered C2a criterion by
two orders of magnitude and a sign (Section~\ref{sec:monopole}); no
computed sector currently supplies $S \simeq 276$.  Were such a
sector found, the smallness of dark energy would be the same kind of
exponential hierarchy that separates $\Lambda_{\mathrm{QCD}}$ from
the Planck
scale.  Identifying the precise topological sector whose susceptibility
sets $\LDE$ is named as the central computation in
Section~\ref{sec:open}.

% ===================================================================
\section{The Planckian floor as a counterterm convention}
\label{sec:cc}
% ===================================================================

The induced gravitational action of \cite{Cheng:2026grav} has, as its
leading vacuum term, a cosmological constant of Planckian size,
\beq
\Gamma_{\mathrm{ind}} \supset c_0\, N\,\Lambda^4
\int\! d^4x\,\sqrt{g},
\qquad N\Lambda^4 \sim \Mpl^4 .
\eeq
Two things must be said about it, one familiar and one decisive for
what follows.

The familiar point is that the coefficient $c_0$ is
\emph{scheme-dependent}---a constant counterterm, exactly as the
induced Newton constant $\Mpl^2$ of \cite{Cheng:2026grav} was found to
be scheme-dependent.  Its renormalized value is fixed by a single
subtraction, conventionally chosen so that the potential minimum sits
at the observed floor.  Why that floor is near zero rather than
Planckian---$\rhoL/\Mpl^4 \sim 10^{-120}$
\cite{Weinberg:1988cp}---is the cosmological-constant problem in its
standard form, and we do not solve it: it is shared with every theory
that contains gravity, and we record it as the one input the
construction does not supply.

The decisive point is that a \emph{constant} counterterm can remove
only the constant part of the vacuum energy.  The
field-\emph{dependent} instanton potential, Eq.~\eqref{eq:Vtheta},
\emph{cannot} be cancelled by it: shifting $c_0$ moves the floor but
leaves the shape $-K\cos(M\theta)$ intact.  A quintessence field
frozen above this minimum therefore carries a positive energy that no
choice of the constant counterterm can subtract away.  This is what
separates the dynamical dark energy developed below from the static
floor: the floor is a convention, the residual is physical.

% ===================================================================
\section{Quintessence phenomenology}
\label{sec:quint}
% ===================================================================

\subsection{The surviving residual is PNGB quintessence}

The surviving field-dependent residual is a small, slowly-varying
energy density---i.e.\ quintessence \cite{Ratra:1987rm,Caldwell:2005tm}.
Equation~\eqref{eq:Vtheta} is a shallow cosine potential for a compact
pseudo-Goldstone; a field of this type with a Planck-scale decay
constant is the pseudo-Nambu--Goldstone (PNGB) quintessence of Frieman,
Hill, Stebbins and Waga \cite{Frieman:1995pm}.  The framework does not
postulate this structure: it is the generic low-energy content of the
spontaneously broken approximate phase symmetry that the substrate
already realizes, now read on the corrected (real/chiral) vacuum rather
than the retired imaginary one.

\subsection{No second field is required}

An earlier version of this programme, in which the angular mode was
read as the \emph{dark-matter} field, faced an over-determination: a
single mode could not carry both an $\mathrm{eV}^4$-scale potential
(the height implied by a $10^{-27}\,\mathrm{eV}$ dark-matter mass) and
the $(\mathrm{meV})^4$ scale of dark energy, and a \emph{second}
pseudo-Goldstone had to be postulated.  The parity analysis of
Section~\ref{sec:axion} removes this requirement.  The matter-coupled,
clustering degree of freedom---the one relevant to galactic
dynamics---is necessarily a \emph{scalar}, not the pseudoscalar
angular mode; in the companion galactic paper it manifests not as a
relic particle but as an acceleration-scale (MOND-like) modification
of the gravitational response (Section~\ref{sec:mond}).  The angular
pseudoscalar is then free to serve as the dark-energy axion
candidate, and a single
complex condensate suffices: its radial/scalar sector governs the
galactic acceleration scale, its angular/pseudoscalar sector the
cosmological one.  The economy the earlier reading lacked is restored.

\subsection{Slow-roll conditions and thawing phenomenology
(conditional)}

\emph{Status note (v6.0).}  With the transmutation mechanism
terminated (Section~\ref{sec:monopole}), no computed sector currently
supplies the potential scale $\LDE^4$; the $R\tilde R$ route is the
sole registered backup.  The phenomenology of this subsection is
therefore \emph{conditional}: it describes what follows \emph{if} a
potential of the protected cosine form arises at the observed scale
from any source.  The parity smoothness theorem and the PNGB
structure are unaffected.

Writing the quintessence field as $\phi$ with
\beq
V(\phi) = \Lambda_{\mathrm{DE}}^4\,
\bigl[\,1 + \cos(\phi/f)\,\bigr],
\eeq
the present dark-energy density requires
$\Lambda_{\mathrm{DE}}^4 \sim \rhoL \sim (2.3\,\mathrm{meV})^4$.  For
the field to act as dark energy today it must be slowly rolling,
i.e.\ frozen by Hubble friction until recently, which requires
\beq
m_\phi^2 \;\simeq\; \frac{\Lambda_{\mathrm{DE}}^4}{f^2}
\;\lesssim\; H_0^2 \;\sim\; \frac{\rhoL}{3\Mpl^2}
\quad\Longrightarrow\quad
f \;\gtrsim\; \Mpl .
\eeq
A super-Planckian decay constant is the standard, and standardly
uncomfortable, requirement of PNGB quintessence; in the present
framework it is a nontrivial constraint on the \emph{single}
condensate's own order parameter, not a free choice: the decay
constant is $f \simeq \sqrt{2}\,v$ (after the flavor-trace
normalization noted in Section~\ref{sec:closed}), and
$v \sim \Lambda/g$ with the lattice scale sitting at or above the
Planck mass in the companion gravitational budget, so
$f \gtrsim \Mpl$ is consistent with---though not yet derived
from---the microscopic parameters.  When such a
field begins to thaw, its equation of state departs upward from
$-1$,
\beq
1 + w_0 \;\simeq\; \frac{\Mpl^2}{3}
\left(\frac{V'}{V}\right)^2_{\!0}
\;\sim\; \left(\frac{\Mpl}{f}\right)^2
\sin^2\!\frac{\phi_0}{f},
\eeq
so the candidate is a \emph{thawing} quintessence with $w \gtrsim -1$
evolving in time, rather than a cosmological constant.

\subsection{What is explained and what remains}

After the C2a termination the ledger reads as follows.  What
survives: the \emph{form}---a protected cosine (PNGB) potential for a
compact pseudo-Goldstone---and the \emph{smoothness}: no clustering,
no fifth force, by parity; these are structural and
measurement-proof.  What is conditional: the \emph{kind}---a
dynamical, thawing component---follows \emph{if} a potential of the
protected form arises at the observed scale from any source.  What is
unexplained: the \emph{amplitude} $\LDE^4$ itself; its one computed
candidate is terminated (Section~\ref{sec:monopole}), and the
equilibrium-vacuum audit (Section~\ref{sec:volovik}) is the
registered successor question.  What does not follow is the near-vanishing of the
vacuum \emph{floor}: the constant counterterm of
Section~\ref{sec:cc} must still be tuned so that the minimum sits near
zero, and this is the standard cosmological-constant problem, untouched
here.  The construction thus separates cleanly into a part it
explains---the dynamical residual---and the one part it shares with all
of physics---the floor.  This is a genuine reduction of the problem,
not a relocation of it: the residual is computed, not fitted.

% ===================================================================
\section{Unification with the galactic sector: the
MOND--$\Lambda$ coincidence}
\label{sec:mond}
% ===================================================================

The scale $\LDE$ that fixes the dark-energy density also fixes a
characteristic \emph{acceleration}.  The natural acceleration built
from the dark-energy scale and the Planck mass is
\beq
a_0 \;\sim\; \frac{\LDE^2}{M_{\mathrm{Pl}}^{\mathrm{(full)}}}
\;\approx\; 2\times10^{-10}\ \mathrm{m\,s^{-2}}
\;\sim\; \frac{cH_0}{2\pi},
\label{eq:a0}
\eeq
within a factor of two of the value
$a_0^{\mathrm{obs}}\approx1.2\times10^{-10}\ \mathrm{m\,s^{-2}}$ at
which galactic rotation curves depart from their Newtonian baryonic
prediction.  In the companion galactic paper \cite{Cheng:2025sparc}
this same scale appears as the crossover of an acceleration-dependent
(k-mouflage) response of the condensate's scalar sector: the force is
screened where the local acceleration exceeds $a_0$ (Solar-System
scales) and active where it falls below (galaxy outskirts), with the
crossover set by $(\partial\chi)^2\sim\rho_{\mathrm{DE}}\sim\LDE^4$.

The numerical coincidence $a_0\sim cH_0\sim c^2\sqrt{\Lambda}$ has been
noted since the inception of MOND and has resisted explanation in
$\Lambda$CDM, where the galactic and cosmological scales are
unrelated.  Here it is structural: the galactic and dark-energy
phenomenologies share a single scale because both are infrared
aspects of one condensate.  Earlier versions of this paper assigned
the galactic response to the condensate's ``scalar sector''; that
assignment does not survive the cross-audit of
Section~\ref{sec:mediator}---the radial mode sits at the cutoff, and
parity forbids the angular mode a linear matter coupling---and the
requisite acceleration-dependent response, if it exists, must instead
be carried \emph{disformally} by the same topological sector that
generates the dark energy.  Eq.~\eqref{eq:a0} fixes only the scale,
which is robust; the mechanism-level unification it suggests is
developed, with its computed constraints, in
Sections~\ref{sec:mediator} and~\ref{sec:monopole}.

% ===================================================================
\section{The mediator gap and the topological funnel}
\label{sec:mediator}
% ===================================================================

\subsection{The gap, stated as a finding}

A cross-audit of the companion papers exposes a structural gap that
this paper is the right place to record, because its resolution runs
through the dark-energy sector.  The galactic phenomenology of
\cite{Cheng:2025sparc} requires a light degree of freedom with a
coherent coupling to unpolarized baryonic matter.  The framework's
documented light modes cannot supply it linearly:
(i) the radial (amplitude) mode sits at the cutoff
($m_\sigma = 2m_{\rm dyn}$, measured);
(ii) the vector channel supports no below-threshold state at the
invariant-universal coupling \cite{Cheng:2026vector};
(iii) the angular mode---this paper's dark-energy axion candidate---is a
pseudoscalar, and the \emph{same parity theorem that guarantees its
smoothness} (Section~\ref{sec:axion}) forbids the linear coupling
$\phi\,T^\mu{}_\mu$ to unpolarized matter; and
(iv) the mixing-induced monopole coupling scales as
$\varepsilon \sim m_\phi^2/\Lambda^2 \sim 10^{-122}$ for a
Hubble-mass field, negligible.
The effective coupling used phenomenologically in
\cite{Cheng:2025sparc} therefore had no assigned microscopic origin.
This is not a fatal inconsistency but an open circuit, and parity
itself dictates the only place it can close: a coupling
\emph{quadratic} in the pseudoscalar,
$\partial_\mu\phi\,\partial_\nu\phi\,T^{\mu\nu}/M^4$---a
\emph{disformal} metric coupling, parity-even, and precisely the
structure employed by relativistic MOND-type theories
\cite{Bekenstein:2004ne}.

\subsection{Two computed negative results close the loop-level
routes}

Whether the framework can generate such a response at the required
strength is a quantitative question, and two computations answer it
for every non-topological route.

\emph{The pairing channel cannot supply the kinetic function.}
For a superfluid phonon, the leading-order effective Lagrangian is
the equation of state, $\mathcal{L}(X) = P(\sqrt{X})$
\cite{Son:2002zn}, and a flat-rotation-curve (MOND) response requires
the non-analytic form $P \propto (\mu-\mu_c)^{3/2}$, i.e.\ a
three-body--dominated polytrope $P \propto n^3$.  A mean-field
computation for the relativistic four-fermion pairing model gives
$P \propto (\mu-\mu_c)^{2}$ at the BEC onset (measured exponent
$2.005 \to 2.000$) and $P \propto \mu^4$ in the BCS regime: a
contact two-body vertex yields the two-body exponent, structurally,
for any coupling strength, cutoff, or fermion mass.  In the adopted
$U(2)\times U(2)$ model the point is moot as well as closed: the
exact pairing decomposition contains no scalar-diquark channel at
all, and the substrate carries no light fermion density at galactic
scales.

\emph{The fermion loop cannot supply the scale.}
Rotating the angular mode into an axial background
$b_\mu = \partial_\mu\phi/2f$ and integrating out the substrate
fermion exactly (constant spacelike $b$; dispersion
$E_\pm^2 = p_\perp^2 + (\sqrt{m_{\rm dyn}^2 + p_\parallel^2} \pm
b)^2$) gives a one-loop effective potential that is an even power
series in $b$ to a verified relative precision of $10^{-11}$, with
$|b|^3$ excluded and the leading nonlinearity entering at
$b \sim m_{\rm dyn}$.  The disformal terms
$b_\mu b_\nu$ in composite kinetic tensors are generated---the
dispersion above is manifestly anisotropic, so the parity gap does
close---but with coefficient $\sim 1/(f^2 m_{\rm dyn}^2) \sim
1/M_{\rm Pl}^4$.  At galactic gradients,
$b \sim \LDE^2/M_{\rm Pl} \sim 2\times10^{-33}\,\mathrm{eV}$ against
$m_{\rm dyn} \sim 10^{28}\,\mathrm{eV}$: the loop-level response is
suppressed by $(b/m_{\rm dyn})^2 \sim 10^{-122}$.

\subsection{The funnel}

The two results together leave exactly one source in the framework
for kinetic or disformal structure that turns on at
$X \sim \LDE^4$ rather than $\Lambda^4$: the topological sector
itself---the same $e^{-S}$ object that generates $V(\phi)$.
Every scale in the loop expansion is Planckian; only the
dimensionally transmuted sector knows about the meV.  The
consequence is a conditional upgrade of the MOND--$\Lambda$
coincidence: \emph{if} the acceleration-dependent galactic response
exists in this framework, it and the dark energy are generated by
the same topological fluctuations, and the two phenomenologies stand
or fall together.  This is a strictly stronger---and therefore more
falsifiable---claim than the shared-scale statement of
Section~\ref{sec:mond}; and it has since been \emph{cashed}: the
funnel's unique address was computed (the monopole gas's gradient
response: non-analytic at the correct scale but of wall type,
$X^{1/2}$, with no source), and the C2a measurement of
Section~\ref{sec:monopole} subsequently removed the diluteness the
address relied on.  The galactic side of the funnel is closed at
leading order; the joint-falsifiability structure did its work in
the killing direction.

% ===================================================================
\section{The C2a measurement: execution and termination record}
\label{sec:monopole}
% ===================================================================

\subsection{What was claimed, and on what it rested}

Versions 5.6--5.8 of this paper reported a passed viability gate: the
programme's Skyrme coefficient, converted through the orientation
layer's exact Skyrme--Maxwell identity to a compact-$U(1)$ coupling
$\beta = 32N\kappa_U \in [5,47]$, appeared to bracket the
monopole-loop action $S \simeq 276$ that the observed dark energy
requires, and to place the layer in its Coulomb phase.  Those claims
rested on the July-8 campaign values $\kappa_U \in [0.054, 0.49]$.
A pre-production pilot subsequently identified those values as
\emph{fit-leakage artifacts}: the box-only diagram set is chirally
incomplete (its zero-momentum amplitude $A_0$ does not vanish), and a
fit basis lacking a constant term leaked $A_0$ into the $k^4$
coefficient, diverging as $1/e^4$ in the external momentum---even the
apparent logarithmic growth toward the chiral limit was mimicked by
the leak.  \textbf{The v5.6--5.8 gate claims, including the
Coulomb-phase co-resolution, are retracted.}

\subsection{The corrected measurement}

The C2a computation (repository \texttt{kappa-c2a}, branch history
public) rebuilt the measurement with pre-registered gates: the full
five-class diagram set (box, triangle, bubble, sunset, contact) with
symbolically derived flavour weights; a chiral-completeness gate
requiring the zero-momentum amplitude to cancel \emph{pointwise},
satisfied at machine precision ($10^{-15}$) and independently
reproduced on a second machine with fresh seeds---a gate that itself
caught and corrected a factor-2 error in the contact prefactor (the
closed fermion loop's flavour trace); the Ward-identity anchor fixing
the vertex normalization ($m\,\Pi_P(0) = -\langle\bar\psi\psi\rangle$,
exact); a plateau protocol ($e \le m/3$, three momenta per mass); an
overdetermined operator projection; a five-point mass scan; two
independent evaluators and two independent integrators; and a stacked
error budget.  The result:
\begin{equation}
\kappa_U \;=\; -0.00149 \pm 0.00001,
\end{equation}
mass-independent across the scan (UV-finite, as expected of the
scheme-independent Skyrme term), with magnitude within $6\%$ of
$1/(64\pi^2)$---squarely at the textbook one-loop
(Gasser--Leutwyler) scale for a chiral fermion loop's $O(p^4)$
coefficients.  The measurement is unremarkable; the mechanism's
requirement was not.

\subsection{The pre-registered verdict}

\begin{quote}
If the pinned $\kappa_U$ places $S_{\rm mono} = 640\,\kappa_U$ outside
$[140, 550]$, the dark-energy chain terminates.\\[2pt]
$S_{\rm mono} = 640 \times (-0.00149) = -0.95 \notin [140,550]$:
\textbf{the chain terminates.}
\end{quote}
Doubly decisive: the magnitude is $\sim\!100\times$ short (no
convention or regulator ambiguity in the audit reaches a factor
$97$), and the sign is negative.  Physically: the substrate's quartic
resistance to orientation twisting is of ordinary one-loop strength,
so monopoles are \emph{cheap} ($S \sim 1$), topological events are
dense rather than exponentially rare, and this sector's vacuum
churns rather than whispers---its susceptibility sits at the cutoff
scale, not the meV scale.  The exponential explanation of dark-energy
smallness is not realized here.

\subsection{Consequences and residual checks}

The gravitational $R\tilde{R}$ route of Section~\ref{sec:closed}
remains the sole registered (uncomputed) transmutation backup.  The
negative sign, if confirmed, additionally implies a wrong-sign induced
Maxwell term for the orientation layer's photon candidate and removes
the one-loop Skyrme stabilizer of the skyrmion (baryon)
identification---consequences carried by the companion working note.
\emph{Both registered residual checks have since been completed}:
the coefficient is derived in closed form,
$\kappa_U = -17/(1152\pi^2) = -0.0014952$ exactly (the measured
value agrees to $0.3\%$; the expansion's $m^2$ and $m^4$ corrections
cancel identically, confirming UV-finiteness analytically), and the
negative sign is confirmed under Pauli--Villars regularization
(three heavy-mass points, converging to the sharp-cutoff value).
The termination is thereby Monte-Carlo--independent:
$S_{\rm mono} = -85/(9\pi^2)$, an exact rational multiple of
$1/\pi^2$, outside $[140,550]$ by construction.

% ===================================================================
\section{The equilibrium vacuum: why the deep sea need not gravitate}
\label{sec:volovik}
% ===================================================================

With the transmutation mechanism terminated, this paper returns the
smallness question to its ontological ground and registers the one
candidate answer consistent with everything measured.  The ontology
(the author's, throughout this programme): the internal-dimension
sector of the condensate is an energy reservoir---a deep sea---of
which the four-dimensional lattice vacuum is the surface.  Hiddenness
alone, however, is not an answer: energy in internal dimensions
gravitates in four dimensions.  The candidate answer is of
\emph{equilibrium} type, due to Volovik \cite{Volovik:2003fe}: for a
\emph{self-sustained} quantum condensate---one that maintains its own
equilibrium without external pressure---the thermodynamic identity
forces the effective vacuum pressure, and with it the gravitating
vacuum energy, to vanish identically in equilibrium.  Not tuned to
zero; zero by the definition of equilibrium.  What gravitates is then
exclusively the \emph{deviation} from equilibrium: waves, knots,
defects, and the small residual by which an expanding universe fails
to be in exact equilibrium.

Two honesty clauses accompany this.  First, the known weakness of the
Volovik programme: the relaxation dynamics that would pin the
present-day residual (and its magnitude, $10^{-122}$ or otherwise)
has never been derived, there or here; the ``thickness of the ice''
remains a price-tagged unknown.  Second, and constructively: in this
framework the statement is \emph{testable at the level of first
principles}, because gravity is induced rather than postulated.
Whether the emergent Einstein equations of the companion gravity
paper source the full condensate energy or only its deviations from
equilibrium is a well-posed audit of the induced source term within
existing machinery.  \textbf{This audit replaces the terminated
transmutation computation as the central open computation of this
paper.}  Its pre-registered stakes: if the induced source contains
the equilibrium ground-state energy, the framework inherits the
cosmological-constant problem in full and the sea ontology loses its
explanatory candidate; if it does not, the smallness question reduces
to the residual's dynamics.

\subsection*{Outlook: expansion as accretion (picture, not derivation)}

The author's cosmological picture, recorded with its status marked:
the lattice grows---expansion is the accretion of new cells, paid
from the sea.  This is kin to sequential-growth dynamics in causal
set theory \cite{Sorkin:2003bx} and to the emergent-space account of
expansion of Padmanabhan \cite{Padmanabhan:2012ik}, in which the
production of spatial degrees of freedom drives, and is driven by,
the Friedmann dynamics.  Within it, the constancy of
$\rho_{\rm DE}$ becomes a bookkeeping identity rather than a
mystery---each new cell carries the same vacuum quota---which is the
legitimate successor of this programme's earliest reading of the
dark energy (``total lattice count''), retired as a mechanism and
revived as an accounting.  The picture owes one equation (the
accretion rate's relation to $H$) before it is physics, and carries
one falsifiable edge worth recording: if accretion is strictly
uniform, cells produced inside bound systems could leave residual
drag signatures constrained by lunar laser ranging and atomic-clock
comparisons.  It is presented as the author's outlook, not as a
derivation.

\section{Observational handle}
\label{sec:obs}
% ===================================================================

The candidate is falsifiable through the equation of state.  A pure
cosmological constant has $w = -1$ exactly; PNGB thawing quintessence
predicts $w_0 > -1$ with $dw/da < 0$ (less negative in the past,
approaching $-1$ at early times)---a definite, distinguishable
signature.  Current baryon-acoustic-oscillation data from DESI
\cite{DESI:2024} favour an evolving equation of state over a constant
$\Lambda$, with a Chevallier--Polarski--Linder fit pulling toward
$w_0 > -1,\ w_a < 0$, i.e.\ the thawing class, at a significance that
has grown but is not yet decisive and that depends on the supernova
sample combined with the BAO data.  We do not over-read this: the
preference may yet be a systematic.  The point of consistency is
nonetheless sharp---and, post-C2a, strictly conditional.  \emph{If}
another microscopic source generates a protected field-dependent
residual at the observed scale, that construction is not free to
predict a constant $\Lambda$: it is by construction dynamical and
falls in the thawing class the data currently prefer.  The
equilibrium-vacuum candidate of Section~\ref{sec:volovik}, by
contrast, makes no thawing commitment yet: whether its residual is
near-constant or dynamical depends on the relaxation/accretion
dynamics that has not been derived.  A confirmed evolving
$w_0 > -1$, $w_a < 0$ would single out the thawing pseudo-Goldstone
class; a strictly constant $w = -1$ would disfavour that class while
leaving the equilibrium-subtraction question untouched.

% ===================================================================
\section{Summary and the central open computation}
\label{sec:open}
% ===================================================================

The dark-energy sector of this framework can be summarized as follows.

\begin{itemize}
\item The condensate contains a parity-protected angular-mode
  \emph{candidate} carrier for dark energy, with the protected
  potential form $V(\phi)=\LDE^4[1+\cos(\phi/f)]$---but no computed
  mechanism currently fixes the amplitude $\LDE^4$
  (Section~\ref{sec:axion}).  The active explanation under test is
  instead equilibrium subtraction of the deep vacuum sea
  (Section~\ref{sec:volovik}).
\item The Planckian vacuum floor is a scheme-dependent counterterm; a
  constant counterterm sets the minimum but \emph{cannot} remove the
  field-dependent residual, which is therefore physical
  (Section~\ref{sec:cc}).  Why the floor is near zero remains the
  standard cosmological-constant problem, the one input not supplied.
\item The field is smooth because parity decouples a pseudoscalar
  from unpolarized matter---a theorem, unaffected by any measurement
  below.  The scale $\LDE^4$ is \emph{not} currently generated by
  anything computed: the transmutation reading survives only as a
  generic possibility (Sections~\ref{sec:axion},
  \ref{sec:monopole}).
\item \emph{Retired historical motivation, recorded for
  provenance}: earlier versions proposed that the same topological
  sector carries both the dark energy and a disformal galactic
  response, upgrading the MOND--$\Lambda$ coincidence to joint
  falsifiability.  That structure was cashed in the double-negative
  direction: the galactic routes are closed at leading order
  (companion galactic paper) and the dark-energy chain is terminated
  (Section~\ref{sec:monopole}).  The scale coincidence
  $a_0\sim\LDE^2/\Mpl\sim cH_0/2\pi$ survives as a relation between
  two observed numbers (Section~\ref{sec:mond}).
\item The transmutation mechanism was constructed, pre-registered,
  and \emph{executed to termination}: the first chirally complete
  measurement of the Skyrme coefficient gives
  $\kappa_U = -0.00149(1)$, ninety-seven times below the requirement
  and of opposite sign; earlier viability claims are retracted as
  fit-leakage artifacts (Section~\ref{sec:monopole}).  The
  smooth-field Hopf-density theorem survives; the gravitational
  $R\tilde R$ route is the sole registered backup.
\item The smallness question is returned to the equilibrium-vacuum
  candidate (Volovik): a self-sustained condensate does not
  gravitate in equilibrium; whether the framework's induced gravity
  realizes this is the new central computation
  (Section~\ref{sec:volovik}).
\item If a protected cosine potential at the observed scale is
  generated by another source, the resulting field belongs to the
  thawing PNGB class ($w_0>-1$ evolving, currently favoured by
  DESI), with the lattice band for $f/\Mpl$ mapping onto the thawing
  rate---conditional phenomenology, not a mechanism's prediction
  (Section~\ref{sec:obs}).
\end{itemize}

The gate ladder of the previous versions was climbed and broke at
its first rung, exactly as its pre-registration allowed: C2a
terminated the chain, and C2b--C2d fall with it.  The paper's open
front is now single and different in kind: the induced-source audit
of Section~\ref{sec:volovik}---does the emergent Einstein equation
weigh the equilibrium sea, or only its deviations?  Its two residual
side-checks (closed-form $\kappa_U$; second-regulator sign) close
the C2a record.  With the audit done---either way---the dark-energy
sector will again rest entirely on computed ground.

% ===================================================================
\section*{Acknowledgments}
The author thanks Claude (Anthropic) for assistance with the
vacuum-structure analysis, the quintessence literature, the
mediator-gap audit and its two negative-result computations, the
pullback theorem and monopole re-scoping of
Section~\ref{sec:monopole}, and manuscript preparation.  The
derivations were developed with substantial AI assistance and the
numerical results of Sections~\ref{sec:mediator}
and~\ref{sec:monopole} were reproduced from archived scripts within
this collaboration; the gate-ladder computations have since been carried out and their
verdicts recorded: C2a (executed, chain-terminating,
Section~\ref{sec:monopole}), C2d (executed in the companion galactic
programme: correct crossover scale, wrong exponent, no source);
C2b--C2c fall with C2a.

% ===================================================================
\begin{thebibliography}{9}

\bibitem{Cheng:2026grav}
Z.~Cheng,
\emph{Emergent gravity from a lattice fermion condensate:
light-cone universality, counterterm degeneracies, and universal
mass-logarithmic coefficients},
companion paper (2026), in preparation.
% NOTE: version numbers removed from bib entries to prevent drift;
% current companion versions are tracked in the header comment.

\bibitem{Cheng:2025sparc}
Z.~Cheng,
\emph{Yukawa-kernel dark matter halos against SPARC rotation
curves}, in preparation (2026).

\bibitem{Cheng:2026vector}
Z.~Cheng,
\emph{The vector channel of a lattice fermion condensate: a massive
composite vector and the sign of the induced Newton constant},
companion paper (2026), in preparation.

\bibitem{Weinberg:1988cp}
S.~Weinberg,
\emph{The cosmological constant problem},
Rev.\ Mod.\ Phys.\ \textbf{61} (1989) 1.

\bibitem{Ratra:1987rm}
B.~Ratra and P.~J.~E.~Peebles,
\emph{Cosmological consequences of a rolling homogeneous scalar
field},
Phys.\ Rev.\ D \textbf{37} (1988) 3406.

\bibitem{Frieman:1995pm}
J.~A.~Frieman, C.~T.~Hill, A.~Stebbins and I.~Waga,
\emph{Cosmology with ultralight pseudo Nambu--Goldstone bosons},
Phys.\ Rev.\ Lett.\ \textbf{75} (1995) 2077
[arXiv:astro-ph/9505060].

\bibitem{Caldwell:2005tm}
R.~R.~Caldwell and E.~V.~Linder,
\emph{The limits of quintessence},
Phys.\ Rev.\ Lett.\ \textbf{95} (2005) 141301
[arXiv:astro-ph/0505494].

\bibitem{DESI:2024}
DESI Collaboration,
\emph{DESI 2024 VI: cosmological constraints from the measurements of
baryon acoustic oscillations} (2024)
[arXiv:2404.03002]; and DESI DR2 cosmological results (2025).

\bibitem{Son:2002zn}
D.~T.~Son,
\emph{Low-energy quantum effective action for relativistic
superfluids},
arXiv:hep-ph/0204199.

\bibitem{Bekenstein:2004ne}
J.~D.~Bekenstein,
\emph{Relativistic gravitation theory for the modified Newtonian
dynamics paradigm},
Phys.\ Rev.\ D \textbf{70} (2004) 083509
[arXiv:astro-ph/0403694].

\bibitem{Polyakov:1976fu}
A.~M.~Polyakov,
\emph{Quark confinement and topology of gauge theories},
Nucl.\ Phys.\ B \textbf{120} (1977) 429.

\bibitem{Volovik:2003fe}
G.~E.~Volovik,
\emph{The Universe in a Helium Droplet},
Int.\ Ser.\ Monogr.\ Phys.\ \textbf{117}, Oxford University Press
(2003).

\bibitem{Padmanabhan:2012ik}
T.~Padmanabhan,
\emph{Emergence and Expansion of Cosmic Space as due to the Quest
for Holographic Equipartition},
arXiv:1206.4916 [hep-th].

\bibitem{Sorkin:2003bx}
R.~D.~Sorkin,
\emph{Causal sets: Discrete gravity},
in \emph{Lectures on Quantum Gravity}, Springer (2005)
[arXiv:gr-qc/0309009].

\bibitem{Guth:1979ik}
A.~H.~Guth,
\emph{Existence proof of a nonconfining phase in four-dimensional
$U(1)$ lattice gauge theory},
Phys.\ Rev.\ D \textbf{21} (1980) 2291.

\bibitem{McGaugh:2016leg}
S.~S.~McGaugh, F.~Lelli and J.~M.~Schombert,
\emph{Radial acceleration relation in rotationally supported
galaxies},
Phys.\ Rev.\ Lett.\ \textbf{117} (2016) 201101
[arXiv:1609.05917].

\bibitem{Lelli:2017vgz}
F.~Lelli, S.~S.~McGaugh, J.~M.~Schombert and M.~S.~Pawlowski,
\emph{One law to rule them all: the radial acceleration relation of
galaxies},
Astrophys.\ J.\ \textbf{836} (2017) 152
[arXiv:1610.08981].

\end{thebibliography}

\end{document}
